% Pursuing Stacks (cotes 134-2 to 134-8): archive cote 134-3, site page 17 (PDF page 18), typescript page 17.
% typescript. Draft compared: ps2.tex:766-819.
% AI-transcribed by Claude Opus 5.5 (claude-opus-5-5, adaptive thinking, effort high), Message Batches API, batch ps1 (msgbatch_01Rv5hCaNTtUt4JHx1yqoHPL), 2026-11-12; not yet checked by eye.
\scanpage{134-3/17}
\mspage{\hand{17}}

appropriate term\xout{s}), which actually need not be topological cells in the
strict sense, but rather conical (and hence contractible) spaces,
each being homeomorphic to the cone \deleted{defined} over the union of all strict-%
\margin{$\times$} ly smaller strata (this union is \uline{compact}). The ordered set associated
to such a (conically) stratified space X is just the set of \inserted{[closed]} strata,
with the inclusion relation, and it can be shown that there is a perfect
dictionary between the topological objects (at least in the case of
finite or locally finite stratifications), and the corresponding (finite
or locally finite) ordered sets, via a ``topological \uline{realization functor}''
\[ K \hand{\longrightarrow} \hand{|}K\hand{|} \]
from ordered sets to (conically stratified) topological spaces. As a
\margin{$\times$} matter of fact, when K is finite, \hand{|}K\hand{|} is endowed with a canonical
triangulation (the so-called barycentric subdivision), the combinato-%
rial model of which \inserted{(or ``maquette'')} is as follows: the vertices are in one-to-one corres-%
pondance with the elements of K (they correspond to the vertices of the
corresponding cones), and the combinatorial simplices \over{\unsure{are}}{are} the \inserted{``flags'', or} subsets
of K which are totally ordered for the induced order. This is still OK
when K is only locally finite (restricting of course to subsets of K
which are finite, when describing simplices), but in any case we can
define (via infinite direct limits\hand{\unsure{§}} in the category (Spaces) of all to-%
pological spaces) the geometric re\over{a}{a}lization of K, \deleted{and} together with
an interpretation of K as the ordered set of strata of \hand{|}K\hand{|} . As a mat-%
ter of fact, w\over{\unsure{ea}}{e} get a canonical isomorphism (nearly tautological)
\[ \hand{|}K \hand{\simeq} \hand{|}\text{Nerve } K\hand{|} \hand{,} \]
where in the second member, we have writ\hand{t}en shortly K for the category
define\over{\unsure{e}}{d} by K. I did not reflect whether it was reasonable to expect
that the categor\over{y}{i}es (Preord) and (Ord) of preordered and ordered sets
are model categories, or even closed model categories, in Quillen's
precise sense - but it is clear though that using ordered sets we'll
get practi\over{\unsure{u}}{c}ally any homotopy type, in any case any homotopy type which
can be described in terms of loca\over{\unsure{l}}{ll}y finite triangulations. However
it should be noted that the inclusion functors into (Cat) or (Ss sets)
\[ \margin{,}\ \text{\over{\unsure{(Pr}}{(O}rd)} \hand{\hookrightarrow} \text{(Pr\hand{\unsure{e}}ord)} \hand{\hookrightarrow} \text{(Cat)} \hand{\hookrightarrow} \text{(Ss sets)}\hand{,} \]
while giving the correct results on homotopy types, do not satisfy \deleted{thex}
Quillen's general conditions on pairs of adjoint functors between model
categories - namely the adjoint functor say
\[ \text{(Cat)} \hand{\longrightarrow} \text{(Preord)}\hand{,} \]
associating to a category A the set Ob A with the obvious order relati-%
on, does \uline{not} commute to formation of homotopy types, as we see in the
trivial case when \deleted{\unsure{thex}} Ob A \inserted{[but not \unsure{Fl}(A)],} is reduced to one point ...
